% !TEX root = ../main.tex
\chapter{Teori og baggrund}
\label{ch:teori}

% Mål: ca. 2.5--3 sider. Skal vise "central anvendt teori" jf. studieordning §4.
% Hold det stramt og relevant for ÉN arm -- ikke et generelt robotik-kursus.

\section{Anvendelseskontekst}
\label{sec:anvendelseskontekst}

Wattson er oprindeligt udviklet til vedligeholds- og
inspektionsopgaver i højspændingsanlæg. Rapporten behandler ikke en
konkret højspændingsanvendelse, men konteksten er taget med, fordi
den begrunder flere af de design- og komponentvalg, integrationen
arbejder ud fra.

Højspændingsanlæg er bygget til mennesker. Betjeningspaneler,
ventiler, afbrydere og gangbroer sidder i menneskehøjde og er
dimensioneret til menneskelig rækkevidde.

En klassisk industrirobot på et fast fundament dækker derfor kun en
lille del af et anlæg, hvis ikke miljøet bygges om. En
menneskeformet robot kan i princippet færdes i den eksisterende
infrastruktur uden tilpasning. Det forklarer valget af en to-armet,
opretstående konstruktion med en rækkevidde tæt på en voksen
persons og en antropomorf arm med 7~DOF.

Anvendelseskonteksten forklarer også flere af de arkitekturvalg,
integrationen arbejder ud fra:

\begin{itemize}
    \item Hver arm har sin egen serielle servobus med en dedikeret
          ESP32-mikrocontroller som bus-master. Delsystemerne kan
          derfor idriftsættes og fejlfindes hver for sig, og én arm
          kan bygges og afprøves uafhængigt af resten af robotten.
    \item Den eksisterende softwarestak understøtter teleoperation
          via VR-udstyr. Det kræver pålidelig og lav-latens
          kommunikation mellem operatør, Jetson og servobus, og det
          stiller krav til buslayout og kabelføring, som gælder
          uanset om armen styres via VR eller via en GUI.
    \item Mekaniske og elektriske grænseflader er designet, så de
          kan demonteres uden specialværktøj, og printede
          slidkomponenter kan udskiftes i drift.
\end{itemize}

% Princip: dette kapitel introducerer KUN den baggrund, som rapportens
% senere kapitler faktisk bygger videre på. Det er ikke et generelt
% kursus i regulering eller robotik.

\section{Seriebus-servoen som integrations-komponent}
% Hvad er en "intelligent" servo (indbygget mikrocontroller, encoder,
% intern regulator). Hvorfor det fra integratorens perspektiv er en sort
% boks man sender positionskommandoer til. Konfiguration via gain_P/I/D
% i JSON-filer nævnes som ydre interface -- ingen PID-teori udledes,
% fordi rapporten ikke tuner regulatoren. Bruges i Kap.~\ref{ch:analyse}
% og Kap.~\ref{ch:idrift}.

\section{Kommunikation: USB-serial og ROS~2}
% Halv-duplex TTL-bus mellem ESP32 og servoer, USB-CDC mellem Jetson og
% ESP32. ROS~2 publish/subscribe-model i det omfang læseren skal kunne
% følge node-diagrammet i Kap.~\ref{ch:analyse} og forstå control_frequency-
% diskussionen i Kap.~\ref{ch:idrift}.

\section{3D-print som konstruktionsmateriale}
% Anisotropi (lag-retning vs. belastningsretning), backlash i printede
% koblinger, kryb under konstant last, slid i bevægelige interfaces.
% Direkte input til materialevalg i Kap.~\ref{ch:mekanik}, til
% observationen af den faldende udstrakte arm i Kap.~\ref{ch:idrift}
% og til diskussionen i Kap.~\ref{ch:diskussion}.
